\documentclass[10pt,a4paper]{amsart}
\usepackage{amsmath,amssymb,amsthm,graphicx}
\usepackage{verbatim}
\newenvironment{reference}{\comment}{\endcomment}
%\newenvironment{reference}{}{}
\newenvironment{slogan}{\comment}{\endcomment}
\newenvironment{history}{\comment}{\endcomment}

% For commutative diagrams we use Xy-pic
\usepackage[all]{xy}

% We use 2cell for 2-commutative diagrams.
\xyoption{2cell}
% Theorem environments.
%
\theoremstyle{plain}
\newtheorem{theorem}[subsection]{Theorem}
\newtheorem{proposition}[subsection]{Proposition}
\newtheorem{lemma}[subsection]{Lemma}

\theoremstyle{definition}
\newtheorem{definition}[subsection]{Definition}
\newtheorem{example}[subsection]{Example}
\newtheorem{exercise}[subsection]{Exercise}
\newtheorem{situation}[subsection]{Situation}

\theoremstyle{remark}
\newtheorem{remark}[subsection]{Remark}
\newtheorem{remarks}[subsection]{Remarks}

\numberwithin{equation}{subsection}

% Macros
%
\def\lim{\mathop{\mathrm{lim}}\nolimits}
\def\colim{\mathop{\mathrm{colim}}\nolimits}
\def\Spec{\mathop{\mathrm{Spec}}}
\def\Hom{\mathop{\mathrm{Hom}}\nolimits}
\def\Ext{\mathop{\mathrm{Ext}}\nolimits}
\def\SheafHom{\mathop{\mathcal{H}\!\mathit{om}}\nolimits}
\def\SheafExt{\mathop{\mathcal{E}\!\mathit{xt}}\nolimits}
\def\Sch{\mathit{Sch}}
\def\Mor{\mathop{\mathrm{Mor}}\nolimits}
\def\Ob{\mathop{\mathrm{Ob}}\nolimits}
\def\Sh{\mathop{\mathit{Sh}}\nolimits}
\def\NL{\mathop{N\!L}\nolimits}
\def\CH{\mathop{\mathrm{CH}}\nolimits}
\def\proetale{{pro\text{-}\acute{e}tale}}
\def\etale{{\acute{e}tale}}
\def\QCoh{\mathit{QCoh}}
\def\Ker{\mathop{\mathrm{Ker}}}
\def\Im{\mathop{\mathrm{Im}}}
\def\Coker{\mathop{\mathrm{Coker}}}
\def\Coim{\mathop{\mathrm{Coim}}}

% Boxtimes
%
\DeclareMathSymbol{\boxtimes}{\mathbin}{AMSa}{"02}

%
% Macros for moduli stacks/spaces
%
\def\QCohstack{\mathcal{QC}\!\mathit{oh}}
\def\Cohstack{\mathcal{C}\!\mathit{oh}}
\def\Spacesstack{\mathcal{S}\!\mathit{paces}}
\def\Quotfunctor{\mathrm{Quot}}
\def\Hilbfunctor{\mathrm{Hilb}}
\def\Curvesstack{\mathcal{C}\!\mathit{urves}}
\def\Polarizedstack{\mathcal{P}\!\mathit{olarized}}
\def\Complexesstack{\mathcal{C}\!\mathit{omplexes}}
% \Pic is the operator that assigns to X its picard group, usage \Pic(X)
% \Picardstack_{X/B} denotes the Picard stack of X over B
% \Picardfunctor_{X/B} denotes the Picard functor of X over B
\def\Pic{\mathop{\mathrm{Pic}}\nolimits}
\def\Picardstack{\mathcal{P}\!\mathit{ic}}
\def\Picardfunctor{\mathrm{Pic}}
\def\Deformationcategory{\mathcal{D}\!\mathit{ef}}
\usepackage{cleveref}
\crefname{originaldefinition}{original definition}{original definitions}
\Crefname{originaldefinition}{Original definition}{Original definitions}
\makeatletter
\newlabel{homology-definition-filtered-complex}{{012L}{source}{}{sourcepointer-1}{}}
\newlabel{homology-definition-filtered-complex@cref}{{[originaldefinition][0][]012L}{[1][0][]source}{}{}{}}
\makeatother
\crefname{originallemma}{original lemma}{original lemmas}
\Crefname{originallemma}{Original lemma}{Original lemmas}
\makeatletter
\newlabel{lemma-K-injective-embedding-filtration}{{0BKI}{source}{}{sourcepointer-2}{}}
\newlabel{lemma-K-injective-embedding-filtration@cref}{{[originallemma][0][]0BKI}{[1][0][]source}{}{}{}}
\makeatother
\crefname{originaltheorem}{original theorem}{original theorems}
\Crefname{originaltheorem}{Original theorem}{Original theorems}
\makeatletter
\newlabel{theorem-K-injective-embedding-grothendieck}{{079P}{source}{}{sourcepointer-3}{}}
\newlabel{theorem-K-injective-embedding-grothendieck@cref}{{[originaltheorem][0][]079P}{[1][0][]source}{}{}{}}
\makeatother
\crefname{originallemma}{original lemma}{original lemmas}
\Crefname{originallemma}{Original lemma}{Original lemmas}
\makeatletter
\newlabel{lemma-derived-products}{{07D9}{source}{}{sourcepointer-4}{}}
\newlabel{lemma-derived-products@cref}{{[originallemma][0][]07D9}{[1][0][]source}{}{}{}}
\makeatother
\title{Realising two inverse systems of derived objects by a bi-filtered complex}
\author{The Stacks Project Authors}
\date{stacks-project:injectives.tex:local-snapshot-20261012; GNU FDL 1.2 or later}
\begin{document}
\maketitle
\noindent\textit{Editorial note.} All delivered mathematical text is the author's own, copied from the pinned Stacks Project chapter file sites-cohomology.tex; no statement or proof step was rewritten, reordered or removed.\par\medskip
\noindent\textit{Editorial note.} The proof is the author's complete argument; the statement is of the form 'In the situation of Lemma ...', so the statement of that lemma (source lines 2601-2614) is delivered as a context span with role=context to keep the card self-contained.\par\medskip
\noindent\textit{Editorial note.} Every labelled reference occurring inside the counted spans is carried in assembly.external\_references as a non-load-bearing navigation pointer (same-chapter ones at their definition site, cross-chapter ones at the invocation inside the counted span), so no load-bearing mathematics is replaced by a stub. No bibliography citation occurs inside the counted spans.\par\medskip
\noindent\textit{Editorial note.} The same-name check for this card scanned every staged Stacks chapter of this repository -- all sources/stacks-*/pinned-original-source.tex files, the card's own chapter excluded -- and the labels found are recorded in dedup.same\_name\_cross\_chapter on the card.\par\medskip
\noindent\textit{Editorial note.} The preamble is the chapter's own retained minimal preamble (sources/stacks-sites-cohomology/retained-preamble.tex), a verbatim excerpt of the native preamble: the Xy-pic setup and the native theorem declarations with the mathematical macros.\par\medskip
\noindent\textit{Allowed standard prerequisites (editorial).}\begin{itemize}
\item Realisation of one inverse system by a filtered complex. Same chapter, Lemma lemma-represent-by-filtered-complex, delivered on this card as a context span. Supplies the situation of the statement and the existence of a filtered complex realising the first inverse system, which the proof applies twice.
\item Products in the derived category. Same chapter, Lemma lemma-derived-products. Supplies that the product of the two complexes realises the object, which is what allows the two filtrations to be combined on a single complex.
\end{itemize}\medskip
\ifdefined\hypertarget\hypertarget{sourcepointer-1}{}\fi
\noindent\textit{Source reference.} Original definition 012L: sources/stacks-homology/pinned-original-source.tex, source lines 6175--6175. Navigation pointer invoked at proof line 2611; the counted span contains the pointer only and the cited result is not reproduced inside it, so the pointer is registered as a non-load-bearing navigation stub in the established form for Stacks material. No bibliography citation is involved.\par\medskip
\ifdefined\hypertarget\hypertarget{sourcepointer-2}{}\fi
\noindent\textit{Source reference.} Original lemma 0BKI: sources/stacks-injectives/pinned-original-source.tex, source lines 2325--2325. Navigation pointer invoked at proof line 2719; the counted span contains the pointer only and the cited result is not reproduced inside it, so the pointer is registered as a non-load-bearing navigation stub in the established form for Stacks material. No bibliography citation is involved.\par\medskip
\ifdefined\hypertarget\hypertarget{sourcepointer-3}{}\fi
\noindent\textit{Source reference.} Original theorem 079P: sources/stacks-injectives/pinned-original-source.tex, source lines 2024--2024. Navigation pointer invoked at proof line 2552; the counted span contains the pointer only and the cited result is not reproduced inside it, so the pointer is registered as a non-load-bearing navigation stub in the established form for Stacks material. No bibliography citation is involved.\par\medskip
\ifdefined\hypertarget\hypertarget{sourcepointer-4}{}\fi
\noindent\textit{Source reference.} Original lemma 07D9: sources/stacks-injectives/pinned-original-source.tex, source lines 2226--2226. Navigation pointer invoked at proof line 2699; the counted span contains the pointer only and the cited result is not reproduced inside it, so the pointer is registered as a non-load-bearing navigation stub in the established form for Stacks material. No bibliography citation is involved.\par\medskip
\begin{lemma}
\label{lemma-represent-by-filtered-complex}
Let $\mathcal{A}$ be a Grothendieck abelian category. Suppose given an object
$E \in D(\mathcal{A})$ and an inverse system $\{E^i\}_{i \in \mathbf{Z}}$
of objects of $D(\mathcal{A})$ over $\mathbf{Z}$ together with
a compatible system of maps $E^i \to E$. Picture:
$$
\ldots \to E^{i + 1} \to E^i \to E^{i - 1} \to \ldots \to E
$$
Then there exists a filtered complex $K^\bullet$ of $\mathcal{A}$
(Homology, Definition \ref{homology-definition-filtered-complex})
such that $K^\bullet$ represents $E$
and $F^iK^\bullet$ represents $E^i$ compatibly with the given maps.
\end{lemma}

\begin{lemma}
\label{lemma-represent-by-filtered-complex-bis}
In the situation of Lemma \ref{lemma-represent-by-filtered-complex}
assume we have a second inverse system $\{(E')^i\}_{i \in \mathbf{Z}}$
and a compatible system of maps $(E')^i \to E$.
Then there exists a bi-filtered complex $K^\bullet$ of $\mathcal{A}$
such that $K^\bullet$ represents $E$, $F^iK^\bullet$ represents $E^i$,
and $(F')^iK^\bullet$ represents $(E')^i$ compatibly with the given maps.
\end{lemma}

\begin{proof}
Using the lemma we can first choose $K^\bullet$ and $F$.
Then we can choose $(K')^\bullet$ and $F'$ which work for
$\{(E')^i\}_{i \in \mathbf{Z}}$ and the maps $(E')^i \to E$.
Using Lemma \ref{lemma-K-injective-embedding-filtration}
we can assume $K^\bullet$ is a K-injective complex.
Then we can choose a map of complexes
$(K')^\bullet \to K^\bullet$ corresponding to
the given identifications
$(K')^\bullet \cong E \cong K^\bullet$.
We can additionally choose a termwise injective
map $(K')^\bullet \to J^\bullet$ with
$J^\bullet$ acyclic and K-injective.
(To do this first map $(K')^\bullet$ to the cone on the identity
and then apply Theorem \ref{theorem-K-injective-embedding-grothendieck}.)
Then $(K')^\bullet \to K^\bullet \times J^\bullet$ and
$K^\bullet \to K^\bullet \times J^\bullet$
are both termwise injective and quasi-isomorphisms
(as the product represents $E$ by Lemma \ref{lemma-derived-products}).
Then we can simply take the images of the filtrations
on $K^\bullet$ and $(K')^\bullet$ under these maps to conclude.
\end{proof}

\end{document}
